% !PRINTED 596
\weberpage{596}
%<<CONT
ist. Die Zahl $\boldsymbol{s}(u)^2$ erzeugt den Körper $\mathfrak{T}_{\mathfrak{m}}$. Es ist außerdem
nach § 122 für jede Zahl $\varrho$ in $\mathfrak{L}$:

\weq{11}{\varrho \equiv \varrho^{p} \quad (\mathrm{mod}\ \mathfrak{P}),}

wenn $\mathfrak{P}$ ein in $\pi$ aufgehendes Primideal in $\mathfrak{L}$ ist, und folglich ist
für jede ganze Zahl $\tau$ in $\mathfrak{T}_{\mathfrak{m}}$, wenn $\mathfrak{P}'$ ein Primideal in $\mathfrak{T}_{\mathfrak{m}}$ ist:

\weq{12}{\tau^{p} \equiv \tau \quad (\mathrm{mod}\ \mathfrak{P}'),}

wenn die Bedingungen (10) und (1) erfüllt sind.

Daraus ergibt sich:

8. \gesperrt{Damit eine in $\Omega$ existierende Primzahl $\pi$ im Körper $\mathfrak{T}_{\mathfrak{m}}$ in Primideale ersten Grades zerfalle, ist notwendig und hinreichend, daß}
\[
\begin{aligned}
\pi &\equiv 1 \quad (\mathrm{mod}\ 2), \\
\pi &\equiv \pm 1 \quad (\mathrm{mod}\ \mathfrak{m}).
\end{aligned}
\]

Von einer gewissen endlichen Anzahl von Primzahlen (die in
den Diskriminanten der auftretenden Gleichungen aufgehen) ist
dabei abgesehen.

\begin{center}
\textbf{\S\,161.\qquad \gesperrt{Zahlgruppen und Idealgruppen.}}
\end{center}

Wie in § 98 und § 106 bezeichne ich mit $O$ die Gesamtheit
der ganzen und gebrochenen Zahlen des quadratischen Körpers $\Omega$,
nach Ausschluß aller Zahlen, die zu einem beliebig angenommenen
$S$ nicht teilerfremd sind. $S$ kann eine rationale Zahl, eine Zahl
in $\Omega$ oder auch ein Ideal sein. Ich will es den \gesperrt{Exkludenten}
nennen. In den drei Abhandlungen \glqq über Zahlgruppen in algebraischen
 Körpern\grqq{}\,\textsuperscript{1)} habe ich gewisse Systeme von Zahlen aus $O$
unter dem Namen \gesperrt{Zahlgruppen} zusammengefaßt, die dadurch
definiert waren, daß das \gesperrt{Produkt} und der \gesperrt{Quotient} je zweier
Zahlen einer solchen Gruppe in derselben Gruppe enthalten waren.
Eine solche Zahlgruppe kann niemals die Zahl Null enthalten,
enthält aber sicher die Zahl 1\,\textsuperscript{2)}.

Im gleichen Sinne wie von Zahlgruppen können wir auch von
\gesperrt{Idealgruppen} reden, die ebenfalls ihren Exkludenten haben
können. Es kommen hierbei natürlich auch gebrochene Ideale
vor, und man bedient sich dabei nicht ohne Nutzen der Darstellungsweise
 durch die \gesperrt{Funktionale} (Bd. II, § 169).

\begingroup
\footnotesize
\noindent\rule{4cm}{0.4pt}\par
\textsuperscript{1)}\, Mathematische Annalen Bd. 48, 49, 50.

\textsuperscript{2)}\, \gesperrt{Fueter} nennt diese Gruppen \glqq \gesperrt{Zahlstrahlen}\grqq{}. Vgl. \gesperrt{Fueter}, Die
Theorie der Zahlstrahlen I, II; \gesperrt{Crelle}, Bd. 130, 132.
\endgroup
